\section{Transformational Possibilities}

\begin{quotex}
Gornahoor would like to welcome Logres who will be posting articles on Fridays for the next several weeks as I take time to pursue other interests. 
\flright{\textsc{Cologero}}
\end{quotex}

The default fundamentalism inherent in so much of contemporary American or even Western Christianity is a cause of some dismay for traditionalists who yearn to look into this tradition for inspiration and support. Watching modern Christians flail their way into movements like the emergent Church\footnote{\url{https://en.wikipedia.org/wiki/Emerging\_church}} or Radical Orthodoxy\footnote{\url{https://en.wikipedia.org/wiki/Radical\_Orthodoxy}} is frustrating or even infuriating. How can so much well-meaning and almost-right energy be spent in ignoring the deeper cultural connections also inherent in the Faith? Any moment, now, one feels that the thrashing may re-ignite the genuine article almost by accident.

Many critics have argued persuasively that Christianity is inherently ``Judaizing'' and doomed to repeat the innate short-circuits infused within the various cultural patterns it assimilated on its path to ascendancy, and they are invariably helped on by some theologians who actually sing paeans of praise.

Specifically, the literalism of the Scriptures as a Holy Book seem to be a fundamental limiting pattern within dialogue on this subject, within and without the Faith. The supposedly perspicacious Writ is quoted against and for us, designedly to limit certain possibilities and force others, usually something compatible with a global universal political settlement.

To these, one might quote Luther, ``The letter is poison.'' (Schrift ist gift). Obviously, this dichotomy set up in the New Testament is a vehicle for higher purpose, and was a major complaint of Jews against Christians, not without reason. To the Biblical narrative, the traditionalist can bring a deeper awareness of unnoticed patterns heretofore ignored. There is, of course, the great solar tradition hidden within the Davidic dynasty, inaugurated when King David opened the inner sanctuary and ate of the shewbread upon the holy mount of peace, a holy site from very ancient times. Troy Southgate's uneven website drew attention to this, in an article by Bill White\footnote{\url{http://www.rosenoire.org/articles/aryan.php}}. It wasn't as if King David was following the letter of the Law in piercing the veil, and he probably didn't have a rope tied to his leg just in case a burial was needed. The order of Melchizedek was another ancient order connected with Salem, and there is also the strange episode of Namaan bowing down in the temple of Rimmon at the advice of Elijah.

Add to this the strange verses in Paul's writings about ``tutelary spirits'' raising the consciousness of nations, baptism for the dead, and the theory of recapitulation of the All through the Kingship-Headship of Christ, as well as parallelisms that exist between many of Christ's saying and the Vedantic\footnote{\url{http://www.biblicalstudies.org.uk/pdf/ijt/30-1_024.pdf}} tradition, as well as the imagery of Revelation, and the Scriptures could begin to look far more interesting. Ask most Christians what Christ meant when he said each little one had a guardian spirit that beholds the face of the Father directly, and you will be met with utter incomprehension or disbelief. We may dismiss their good opinion and proceed. Yet we need not concede them control over the vehicles, since to do so strips our effort of needful road signs or tools.

For example, the Christian Reconstructionist James Jordan\footnote{\url{http://www.biblicalhorizons.com/pdf/jjne.pdf}} has written an obscure but useful book about the hidden symbols of the Bible that are in plain sight, and Jonathan Edwards (had he lived) might have greatly developed his portrayal of the Platonic\footnote{\url{http://www.jesociety.org/2011/02/10/was-jonathan-edwards-a-gnostic/}} imagery of nature\footnote{\url{http://www.ts.mu.edu/content/65/65.1/65.1.3.pdf}}. Is it an accident that men like Marsilio Ficino pop up within Western streams of thought and religion, \emph{or is it wholly for our sake, for the Now, that they wrote what they wrote}? \emph{To cleanse the doors of perception for the I which aspires in the Now for truth, but has lost its way}? In the dearth of a true tradition, Christians (and those who would use the vehicle as they may) are compelled to re-forge the sword which was broken -- we cannot look to the official institutions themselves for confirmation, aid, or blessing.

To those false sons within the pale, and those who scorn the true Church as utterly bereft of real grace, we ought to reply with Isaac Penington, and apply this unction to the dying Faith:

\begin{quotex}
Every truth is a shadow except the last. But every truth is substance in its own place, though it be but a shadow in another. And the shadow is a true shadow, as the substance is a true substance.
\end{quotex}


These fragments of truth which she has grasped, like some \emph{de copia} of Erasmus, she has shored against her ``ruin''. The work will begin. We will make our inheritance our own, and the letter shall be burned.


\flrightit{Posted on 2011-04-08 by Logres}

\begin{center}* * *\end{center}

\begin{footnotesize}
\begin{sffamily}
\texttt{Spock on 2011-04-08 at 09:13 said:}

Biblical passages don't prove much of anything, only that the scribes or forgers could copy from other manuscripts.

\hfill

\texttt{Cologero on 2011-04-08 at 12:56 said:}

There are some defects in the James Jordan text, particularly in his understanding of the ``creation'' of the world. He takes the position that God could have created any possible world, where by ``possible'' he means any world that James Jordan can conceive in his imagination. However, there is no reason to suppose that whatever James Jordan can imagine is actually realizable. That may be the Jewish interpretation, but it is not the Christian.

Christian understanding states that everything was created through the Logos: ``All things were made by him Logos; and without him was not any thing made that was made.''

So when James Jordan wonders why the world is the way it is, that is because it conforms to a cosmic order, a topic we have addressed many times. As far as Mr. Jordan's objection to Plato, we can assure him that it was John himself who made the connection to the Logos of Greek philosophy. The Alexandrians and Augustine just noted it and drew out the consequences.

\hfill

\texttt{logres on 2011-04-08 at 20:58 said:}

Thanks for pointing that out by articulating it clearly; that's the problem with many branches of Christian tradition, and one I know very personally -- one has a naive faith, then it is ``lost'', and/or one comes upon a writer like Jordan who does dip into the pool a little bit, leaving one with the general idea that all has been said. Whereupon, one shuts one's self off from any further development. In Jordan's case, he's just more interesting than the average theologian, and therefore more helpful and dangerous all at once. Such debates do represent a possible apologetic angle for traditionalists, as Cologero has outlined. And thanks for the welcome.

\hfill

\texttt{Ismo Meinander on 2011-04-09 at 05:29 said:}

Cologero said: Christian understanding states that everything was created through the Logos: ``All things were made by him Logos; and without him was not any thing made that was made.''

There is a very interesting passage in Wolfgang Smith's `The Quantum Enigma' concerning this biblical quotation, in which he discusses the issue of causation and vertical causality in the context of art and creativity. I think it is best for me to quote him straight from the book:

``It emerges that there exists a primary causality which acts, not in some distant past, but in every here and now without exception. All things existing in space and time are not only brought into being, but held in existence, byt this primary causation which derives from a single and indivisible Act. Unlike the kinds of causality with which modern science is concerned -- which may be termed temporal or natural causation -- this primary causality does not act from past to futur by way of a temporary process, but acts directly, unmediated by any chain of temporal events. The question arises now whether this `temporally unmediated' mode of action -- which we shall designate by the adjective vertical -- is the exclusive prerogative of primary causation, or whether perhaps there exist secondary modes of vertical causality. In answer to this question it can be said that the causation effected by an intelligent agent is perforce vertical. Take the case of art in the primitive sence of human making: the entire process hinges in fact upon such a vertical act. What stands at issue in authentic art is a veritable imitatio Dei: the human artist `participates' to some degree in the creative process of the First Cause: `All things were made by Him, and without Him not anything was made' (John 1: 13). But does this mean that all production -- even the shoddiest artifact -- is to be ascribed indisriminately to God Himself? Assuredly not. It is interesting to note, in this connection, that according to the punctuation which became generally accepted in post-medievval times, John 1: 13 actually reads: `All things were made by Him, and without Him was not anything made that was truly made.'' We may take it that the quod factum est refers to what is truly made, and therefore what truly is. The difference, Scholastically speaking, lies in the presence or absence of form: it is form -- a transcendent element! -- that bestows being. Now, the bestowal of form constitues an incurably vertical act of causation.''

So I take it that `vertical causation' is another way of expressing the `eternal now' and in reference to the `creation' it can be said that the world is created by the immanently transcendent divinity anew from moment to moment. So the passage could go also `everything is created by Him and without Him nothing is made that is truly made. I think this conception of `an eternal creation' should be emphasized and elaborted every time when one speaks for example with scientists about `the creation of the world' and the world process: The Primary Act did not act in `the big bang' and then evade itself in some transcendent dimension of Being like the Deists seem to think, but it acts in everything, especially and in potentia conscously through humans who have the possibilities to conceive it unlike animals and non-human Nature that follows its ways instinctually, and nothing could not exist without this immanent participation into the higher states of being either consciously or un-consciously. In the realm of art and human creativity it can be said that true art in some way or another reflects harmoniously the intellect of the Logos.

\hfill

\texttt{Ismo Meinander on 2011-04-09 at 05:40 said:}

Correction: `... and nothing could not exist... ' -- `nothing could exist'. My english is not perfect \& I apologize some possible mistakes.

\hfill

\texttt{Cologero on 2011-04-09 at 19:13 said:}

Thanks for pointing out that passage from Wolfgang Smith. It ties in, too, with Mouravieff's New Age, the era of the Artist ... as you and Prof Smith use the term. Artistic intuition needs to be expressed, and artistic expression is a continuing creative process, not a once and for all times event.

\hfill

\texttt{Cologero on 2011-04-09 at 19:32 said:}

By attempting to recover the symbolic, Jordan is accidentally interesting but he ruins it with his literalness and fundamentalism. As an example, we can point to his misunderstanding of the cosmology of the four elements as a sort of proto-chemistry in the modern sense. Of course not, but as a qualitative understanding of the world, it is adequate and accurate. Could God have created a world without one of the four qualitative elements? Perhaps a world without the liquid state (water)? Inconceivable. Or without energy (fire)? Not at all. God is arbitrary for Jordan, not consubstantial with the Logos.

\hfill

\texttt{logres on 2011-04-10 at 00:49 said:}

Yes, I hadn't thought deeply enough about the why of Jordan being wrong, but ``accidentally interesting'' probably aptly characterizes his work. I read a Catholic writer who thought that Duns Scotus' emphasis on God's will lead to problems in the West, but as I've come across traditionalist writers, I see that even this doesn't get deep enough. The examples are endless: Richard Weaver blaming William of Occam in Ideas Have Consequences, etc., etc. A lot of blaming in the West, not much substance, or rather, not enough to warrant a prophetic tone. Harold Weatherby has an interesting book called ``The Keen Delight'' on the effect of bad philosophy on poetry. He starts with Dante, goes through Hopkins, and ends with Eliot. Of course, his answer is more ``scholasticism'' of the Aquinas-type.

\hfill

\end{sffamily}
\end{footnotesize}

